\section{Problemas}

\begin{frame}{Problemas}
  \begin{itemize}
    \item \textbf{Igualdad de caminos en un árbol.}
      Un caracter en cada arista, y consultas \((u_1,v_1,u_2,v_2)\).
      La string de un camino es la concatenación de las aristas, de \(u\) a \(v\).
      Decir si los dos caminos forman la misma string.

    \item \textbf{Comparación de caminos en un árbol.}
      Las mismas consultas.
      Si las strings son distintas, decir cuál es menor.

    \item \textbf{Árboles isomorfos.}
      Dos árboles. Decidir si son isomorfos.
      Se puede resolver sin hashing.
      Con sum hashing, la implementación queda más corta.

    \item \textbf{ABC 481 F.}
      Grafo dirigido, un dígito del 1 al 9 en cada arista.
      El costo de un camino es el número que forman esos dígitos.
      Para cada nodo, el menor costo desde el nodo 1, módulo \(998244353\).

      {\footnotesize \url{https://atcoder.jp/contests/abc481/tasks/abc481_f}}
  \end{itemize}
\end{frame}

\begin{frame}{Problemas}
  \begin{itemize}
    \item \textbf{TAP 2019 H, Rummaging.}
      Un texto, y varias frases del mismo largo.
      Contar cuántas veces aparecen en el texto, permitiendo a lo sumo un caracter distinto.

      {\footnotesize \url{https://codeforces.com/group/YjFmW2O15Q/contest/102363/problem/H}}

    \item \textbf{TC ARG 2024, día 7, F.}
      \(N\) strings del mismo largo, con caracteres A, B, C y D.
      Hay una sola cuya cantidad de posiciones distintas con cada una de las demás es \(K\).
      Encontrarla.

      {\footnotesize \url{https://codeforces.com/gym/104686/problem/F}}
  \end{itemize}
\end{frame}

\begin{frame}[fragile]{Problemas}
  \begin{itemize}
    \item \textbf{OIA 2022, Curry o Curro.}
      Un diccionario de palabras, cada una con un significado entero.
      Sumar, sobre los pares del mismo largo que difieren en exactamente una letra, la diferencia absoluta de sus significados.

      {\footnotesize \url{https://juez.oia.unsam.edu.ar/#/task/curry/statement}}

    \item \textbf{ABC 250 E.}
      Dos arreglos \(A\) y \(B\).
      Para cada consulta \((x, y)\), decir si el conjunto de valores del prefijo de \(A\) de largo \(x\) es igual al del prefijo de \(B\) de largo \(y\).

      {\footnotesize \url{https://atcoder.jp/contests/abc250/tasks/abc250_e}}

    \item \textbf{Codeforces 1622 F.}
      Un conjunto es cuadrático si el producto de los factoriales de sus elementos es un cuadrado perfecto.
      Dado \(n\), encontrar un subconjunto cuadrático de \(\{1,\ldots,n\}\) de tamaño máximo.

      {\footnotesize \url{https://codeforces.com/contest/1622/problem/F}}

    \item \textbf{Codeforces 1175 F.}
      Un subarreglo es una subpermutación si contiene cada entero desde \(1\) hasta su largo, exactamente una vez.
      Contar cuántas hay.

      {\footnotesize \url{https://codeforces.com/problemset/problem/1175/F}}
  \end{itemize}
\end{frame}

\begin{frame}{Problemas}
  \begin{itemize}
    \item \textbf{Codeforces 763 D.}
      Un árbol.
      Elegir una raíz que maximice la cantidad de subárboles distintos, salvo isomorfismo.
      Alcanza con devolver cualquier raíz óptima.

      {\footnotesize \url{https://codeforces.com/contest/763/problem/D}}
  \end{itemize}
\end{frame}

\begin{frame}[fragile]{Material adicional}
  \begin{itemize}
    \item \textbf{XOR hashing.}
      Tutorial.
      A cada valor se le asigna un número aleatorio, y el XOR reemplaza al conjunto.
      Con eso se puede chequear si un subarreglo es una permutación.

      {\footnotesize \url{https://codeforces.com/blog/entry/85900}}

    \item \textbf{Probabilidad de colisiones.}
      El comentario acota la probabilidad de un falso positivo al chequear una permutación con el XOR de hashes aleatorios de \(B\) bits.
      Esa probabilidad es \(2^{-B}\).
      Con \(B = 64\), es menor que \(10^{-19}\).

      {\footnotesize \url{https://codeforces.com/blog/entry/121682?#comment-1080263}}
  \end{itemize}
\end{frame}
